\section{Background}
\label{sec:background}

SMT solvers~\cite{smtchapter} decide the satisfiability of a first-order logical
formula. They are widely used in program verification, automated reasoning, and
other areas of computer science. SMT solvers (such as Sundance) work by
implementing a Boolean satisfiability (SAT) engine that searches for a
satisfying assignment to the propositional skeleton of the input formula, and a
theory solver that checks the validity of the theory-specific constraints, such
as integer and real arithmetic, algebraic datatypes, and quantifiers. We refer
you to \citet{smtchapter} for a more detailed introduction to SMT solving.

\subsection{Running example}

To make SMT solving and the role of e-matching concrete, consider the
SMT-LIB~\cite{smtlibstandard} query in Figure~\ref{fig:running-example}. The
query declares two functions over an uninterpreted sort \smt{U}: \smt{g}
(binary, line~2) and \smt{h} (unary, line~3). It asserts a top-level disjunction
(line~8); three ground facts (lines~9--11); and a single universally quantified
clause (lines~12--13) with the nested trigger \smt{(g (h x) x)}. The query is
unsatisfiable, and a typical SMT solver discharges it as follows.

The SAT engine first decides the disjunction by case-splitting on its two
disjuncts.

\paragraph{Branch \smt{a = b}.} Asserting \smt{a = b} merges the e-classes of
\smt{a} and \smt{b}. Congruence closure then propagates this merge through the
existing \smt{g}-applications: the term \smt{(g (h a) b)} becomes congruent to
\smt{(g (h a) a)}, so the e-graph now contains an application of \smt{g} whose
first argument is \smt{(h a)} and whose second argument is in the e-class of
\smt{a}. The trigger \smt{(g (h x) x)} therefore matches with \smt{x = a}, and
the quantifier instantiates to \smt{(not (g (h a) a))}. This new lemma
contradicts the term \smt{(g (h a) a)} that congruence closure just derived in
the e-graph, producing a theory conflict.  Notice that we have reached a 
contradiction, but only in the current branch of the search.
The solver records a conflict clause
and backtracks to before the \smt{a = b} decision.

\paragraph{Branch \smt{a = c}.} This branch is similar. Merging the e-classes
of \smt{a} and \smt{c} makes \smt{(g (h a) c)} congruent to \smt{(g (h a) a)}.
The same trigger \smt{(g (h x) x)} matches with \smt{x = a}, the quantifier
instantiates to \smt{(not (g (h a) a))}, and we again obtain a theory conflict.

With both disjuncts refuted, the SAT engine has no remaining assignments to the
disjunction and the solver returns \texttt{unsat}. 

% The ground fact \smt{(g a a)}
% never participates in either conflict --- it is a decoy present only to inflate
% the \smt{g}-table and stress the e-matching algorithm.

\begin{figure}
\begin{smtlib}
(declare-sort U 0) 
(declare-fun g (U U) Bool) 
(declare-fun h (U) U)
(declare-const a U) 
(declare-const b U) 
(declare-const c U)

(assert (or (= a b) (= a c)))
(assert (g (h a) b)) 
(assert (g (h a) c)) 
(assert (g a a))
(assert (forall (x U) (not (g (h x) x)) 
          :pattern ((g (h x) x))))
\end{smtlib}
\caption{A simple SMT query exercising backtracking (via the disjunction) and a
nested trigger pattern. Top-down e-matching scans all \smt{g}-entries;
relational e-matching can start  from the small \smt{h}-table and index-joins into
\smt{g}, skipping the \smt{(g a a)}  term entirely.}
\label{fig:running-example}
\end{figure}

\subsubsection{Takeaways}
There are two key features that are important to highlight here. First, the disjunction
forces the SAT engine to explore two case splits. After choosing one disjunct,
instantiating quantifiers, and propagating consequences, the solver may discover
a conflict and \emph{backtrack} to try the other disjunct. Thus, all of the
solvers datastructures must be rolled back to the state they were in before the first
disjunct.

% Any e-matching index
% built up under one assignment must therefore be undone (or otherwise rolled
% back) when the solver pops a decision level.
Second, the trigger \smt{g(h(x),
x)} is a nested pattern: matching it requires finding a \smt{g}-application
whose first argument is an e-class that contains an \smt{h}-application and
whose second argument is in that same e-class. The two e-matching algorithms
below take very different approaches to finding such matches.

\subsection{E-matching}

\subsubsection{Top-down e-matching}
\label{sec:bg:topdown}

Top-down e-matching~\cite{simplify}, the classical algorithm used in most SMT
solvers, walks the pattern from the top and recursively matches against the
e-graph. For our trigger \smt{(g (h x) x)}, the algorithm enumerates every e-node
of the form \smt{(g t1 t2)} and checks whether
the e-class of \smt{t1} contains an \smt{h}-application \smt{(h s)}, and whether
the e-class of \smt{s} equals the e-class of \smt{t2}. In our running example
after the merge \smt{a = b}, the \smt{g}-table contains the entries \smt{(g (h
a) b)}, \smt{(g (h a) c)}, and \smt{(g a a)}; the algorithm probes all three
even though only the first yields a valid match. More generally, with $n$
entries in the \smt{g}-table and $k$ entries in the \smt{h}-table, top-down
e-matching costs $O(n + k)$ work to locate the $k$ relevant candidates, regardless
of how small $k$ is.

\subsubsection{Relational e-matching}
\label{sec:bg:relational}

Relational e-matching~\cite{relationalematching} treats each function symbol's
e-node table as a database relation and compiles the trigger into a multi-way
join. For \smt{(g (h x) x)}, the trigger is encoded as the conjunctive query
%
\[
  h(x) = v_0 \;\;\wedge\;\; g(v_0,\; x) = root
\]
%

The a query scheduler can choose to start from either atom. Starting from the
\smt{g}-atom means the algorithm must scan all $n$ entries in the \smt{g}-table.
Since there are fewer entries in the \smt{h}-table, a better choice is to start
from the \smt{h}-atom: the algorithm looks up the $k$ entries in the
\smt{h}-table, and for each one probes the \smt{g}-table for entries whose first
argument is the e-class of that \smt{h}-entry. This costs $O(k \log n)$ work,
which is a substantial improvement over top-down e-matching when $k$ is small
relative to $n$.